\section{Unifying framework: a sheaf of CoHAs on the $(2,1)$-stack of equivariance strata}
\label{sec:gluing:unifying}
\label{gluing:part:x-unifying}
The local-to-global CoHA assemblies of \S\S\ref{sec:gluing:abstract}--\S\ref{sec:gluing:obstructions} ---
chart-wise on toric fans, orbifold-inertia-twisted on McKay quotients,
tilting-mediated on derived Morita classes, factorisation-algebra-
structured on Ran spaces, wall-crossing-conjugated on stability
chambers, Borcherds-lifted on lattice-polarised period domains ---
consolidate into one object: a sheaf of CoHAs on a $(2,1)$-stack
$\Stk_{\Str}$ whose points parametrise equivariance strata, whose
$1$-morphisms are stratum-inclusions, whose $2$-morphisms are the
coherences between stratum-compatibility witnesses. The $K3\times E$
CoHA is the deepest stalk of this sheaf.

The four strata introduced in \S\ref{sec:gluing:abstract} are not a partition: a CY$_3$
lives in several at once, and the intersections proliferate. The
subsets $S\subseteq\{i,ii,iii,iv\}$ number $2^{4}=16$. The empty
subset $S=\varnothing$ is excluded on structural grounds: the
Davison--Meinhardt critical CoHA requires a $\Gm$-scaling of
$\Omega_{X}$ for the Joyce--Song sign
$\varepsilon(E,F)=(-1)^{\chi(E,F)}$ to be well-defined
(Joyce--Song \texttt{arXiv:0810.5645}~Def.~5.2 introduces
$\varepsilon$; Thm.~5.14 proves the integration map is an algebra
homomorphism using the $\Gm$-scaling compatibility of $\varepsilon$;
Kontsevich--Soibelman \texttt{arXiv:0811.2435}~\S 6.3 gives the mixed
Hodge module incarnation on the vanishing-cycle sheaf). Every CY
admitting such a scaling lies in at least one of the four strata:
the $\Gm$-scaling is $T^{d}$-diagonal in stratum~(i), the reduced
$\Gm\subset\Aut_{s}(X)$ in stratum~(ii), induced from the ambient
$\Gm$ on $\C^{d}$ before $G$-quotient in stratum~(iii), induced from
the period-domain $\Gm$-action in stratum~(iv). The remaining fifteen
subsets each carry CoHA content. They organise
into a partially-ordered stack under stratum-inclusion, and the
gluing cocycle of $\cH(X)$ depends only on $\Str(X)\subseteq\{i,ii,iii,iv\}$.
Section~10.1 states the classification; \S 10.2 upgrades it to a
sheaf on a $(2,1)$-stack; \S 10.3 binds the construction across the
three volumes through the two-stage factorisation
$\Phi_{3}=\mathrm{Sp}^{\mathrm{ch}}_{\Sigma_{2},C}\circ\Phi^{FA}_{3}$;
\S 10.4 names the frontiers.

\subsubsection{\S 10.1. The fifteen-cell classification}
\label{gluing:sec:10-1-fifteen-cell}

Recall the four strata of \S\ref{sec:gluing:abstract}. Write $\Str_{i}$ for the stratum of
toric $T^{d}$-charts, $\Str_{ii}$ for the reduced
$\Gm\times\Aut_{s}(X)$-stratum (volume-preserving residual
automorphism), $\Str_{iii}$ for the orbifold-inertia stratum
$I(X/G)$, $\Str_{iv}$ for the lattice-polarised period-domain
stratum. For a CY$_{d}$ variety $X$ (or CY$_{d}$ category $\cC$)
write $\Str(X)\subset\{i,ii,iii,iv\}$ for the set of strata $X$
belongs to.

\begin{theorem}[Fifteen-cell classification, structural statement]
\label{gluing:thm:fifteen-cell-classification}
\ClaimStatusConjectured\ \emph{(Theorem-level for singleton cells
$1$--$4$ per Davison--Meinhardt \texttt{arXiv:1601.02479}, Mukai
\texttt{arXiv:math/0408129}, BKR \texttt{arXiv:math/9908027},
Borcherds \texttt{arXiv:alg-geom/9609022}; conjectural in full
generality for pairwise and higher cells.)} For every CY datum
$X\in\Stk_{\Str}$ admitting the Joyce--Song $\Gm$-scaling
(\texttt{arXiv:0810.5645} Def.~5.2), $\Str(X)\neq\varnothing$. The
global CoHA $\cH(X)$ admits a local-to-global presentation whose
gluing cocycle class
\[
[c_{X}]\;\in\;H^{1}\bigl(X,\,\underline{F}_{\Str(X)}\bigr)
\]
takes values in a sheaf of groups $\underline{F}_{\Str(X)}$ depending
only on $\Str(X)\subseteq\{i,ii,iii,iv\}$, according to the
fifteen-cell table below. On singleton cells $\underline{F}$ is one
of the four basic sheaves
\[
\underline{\Pic}_{T^{d}},\qquad
\underline{\Aut_{s}(X)},\qquad
\underline{G},\qquad
\underline{\Gamma_{\mathrm{arith}}};
\]
on higher cells $\underline{F}$ is the fibre product when the two
strata share only the diagonal $\Gm$-scaling of $\Omega_{X}$, and a
semidirect product when one stratum's group acts by outer
automorphism on the other's cocycle sheaf (e.g.\ $\Aut_{s}(X)$ acts
on the conjugacy-class indexing of $G \subset M_{24}$ in cell~$8$,
producing $\underline{\Aut_{s}}\rtimes\underline{G}$).
\end{theorem}

\noindent We enumerate the cells. Each cell is labelled by its stratum
set; the associated cohomology group names the first-Cech cohomology
of the relevant sheaf on the natural site; the representative $X$ is
the simplest Vol~III witness.

\begin{center}
\begin{tabular}{|c|c|c|c|}
\hline
Cell & $\Str(X)$ & Cocycle group & Representative $X$ \\
\hline
1 & $\{i\}$ & $H^{1}(X,\underline{\Pic}_{T^{d}})$ & $\C^{3}$, conifold, banana, SPP \\
2 & $\{ii\}$ & $H^{1}(X/\Aut_{s}(X),\underline{\Aut_{s}(X)})$ & abelian 3-fold $A^{3}$; K3 with Nikulin involution (at $d=2$) \\
3 & $\{iii\}$ & $H^{1}(I(X/G),\underline{G})$ & $[\C^{3}/G]$ with $G\subset SU(3)$ finite \\
4 & $\{iv\}$ & $H^{1}(\cD_{L}/\Gamma_{\mathrm{arith}},\underline{\Gamma_{\mathrm{arith}}})$ & K3 lattice-polarised at a Humbert divisor \\
\hline
5 & $\{i,ii\}$ & $H^{1}(X,\underline{\Pic}_{T^{d}}\rtimes\underline{\Aut_{s}(X)})$ & toric CY$_{3}$ with non-trivial $\Aut_{s}$ (e.g.\ conifold $\rtimes\Z/2$) \\
6 & $\{i,iii\}$ & $H^{1}(I(X/G),\underline{\Pic}\rtimes\underline{G})$ & local $\mathbb{P}^{2}=[\C^{3}/\Z_{3}]^{\mathrm{res}}$ \\
7 & $\{i,iv\}$ & $H^{1}(X,\underline{\Pic}_{T^{d}}\times_{\Gm}\underline{\Gamma_{\mathrm{arith}}})$ & toric degeneration of a $\cD_{L}$-lifted CY at $\cH_{D}$ \\
8 & $\{ii,iii\}$ & $H^{1}(I(X/G)/\Aut_{s},\underline{\Aut_{s}}\rtimes\underline{G})$ & Mathieu $M_{24}\curvearrowright K3$ with residual $\Aut_{s}$ \\
9 & $\{ii,iv\}$ & $H^{1}(\cD_{L}/\Aut_{s},\underline{\Aut_{s}}\rtimes\underline{\Gamma_{\mathrm{arith}}})$ & K3 lattice-polarised at Humbert $\cH_{D}$ \\
10 & $\{iii,iv\}$ & $H^{1}(\cD_{L}/G,\underline{G}\rtimes\underline{\Gamma_{\mathrm{arith}}})$ & Mathieu K3 at Humbert divisor \\
\hline
11 & $\{i,ii,iii\}$ & fibre-product at three strata & $[\C^{3}/\Z_{3}]^{\mathrm{res}}$ with Hesse $\Aut_{s}$ \\
12 & $\{i,ii,iv\}$ & fibre-product at three strata & toric abelian at Humbert \\
13 & $\{i,iii,iv\}$ & fibre-product at three strata & local $\mathbb{P}^{2}$ at Humbert \\
14 & $\{ii,iii,iv\}$ & fibre-product at three strata & $K3\times E$ (ii,iii,iv) \\
\hline
15 & $\{i,ii,iii,iv\}$ & nested fibre-product at four strata & $K3\times E$ with elliptic-fibration toric germs \\
\hline
\end{tabular}
\end{center}

\noindent The key observation: for cells 1--4 the cocycle group is
exactly one $H^{1}$ of one sheaf on one site. For cells 5--10 the
cocycle is a fibre product over the boundary of the two strata. For
cells 11--14 the fibre product nests threefold. For cell 15, where
all four strata meet, the cocycle is a fourfold nested fibre product,
and the canonical witness is $K3\times E$ equipped with its full
four-stratum structure (elliptic fibration giving a partial toric
input to stratum~$(i)$, $\Aut_{s}$ giving $(ii)$, $M_{24}$ giving
$(iii)$, $\Delta_{5}$ giving $(iv)$). Cell~$15$ is the deepest.

\paragraph{Elementary derivation of the singleton cells.}

\emph{Cell 1 (toric alone).} A toric CY$_{3}$ $X$ has a fan $\Sigma$
with top-dimensional cones $\{\sigma_{\alpha}\}$ and charts
$U_{\alpha}\cong\C^{3}$; adjacent cones share a codimension-one face
with overlap $U_{\alpha\beta}\cong\Gm\times\C^{2}$. Each chart carries
the Jordan-triple potential $W_{\alpha}=\mathrm{tr}(x_{\alpha}[y_{\alpha},
z_{\alpha}])$ (Kontsevich--Soibelman~\texttt{arXiv:0811.2435}~\S 6.1);
on $U_{\alpha\beta}$ the two potentials $W_{\alpha}, W_{\beta}$ differ
by a Fourier--Mukai-type kernel
$K_{\alpha\beta}=\cO_{\Delta}\otimes\pi_{1}^{*}\cO(k_{\alpha\beta})$
with $k_{\alpha\beta}\in\Z$ dictated by the primitive ray of the shared
edge (Van den Bergh \texttt{arXiv:math/0211064}~\S 3; Szendroi
\texttt{arXiv:0705.3419} for the conifold two-chart case). Assembled
over the $1$-skeleton of the fan, the $\{k_{\alpha\beta}\}$ form a
$\Cech$ $1$-cocycle with values in the constant sheaf
$\underline{\Pic}_{T^{3}}$ on the nerve of the cover; its class in
$H^{1}(X,\underline{\Pic}_{T^{3}})$ is the cell-$1$ cocycle
(Davison--Meinhardt \texttt{arXiv:1601.02479}~\S 2).

\emph{Cell 2 (reduced $\Aut_{s}$ alone).} The singleton cell~$2$
requires a CY$_{d}$ with non-trivial residual volume-preserving
automorphism group $\Aut_{s}(X)$ and no toric, orbifold, or
lattice-polarised structure. A generic K3 has trivial $\Aut_{s}^{0}$
(Huybrechts, \emph{Lectures on K3}, Ch.~15) and is therefore \emph{not}
in cell~$2$; the canonical $d = 3$ representative is an abelian
3-fold $A^{3} = \C^{3}/\Lambda$ whose $\Aut_{s}^{0} = A^{3}$ is the
translation action (itself $d$-dimensional). A K3 with a single
Nikulin involution $\iota \in \Aut_{s}(K3)$ realises cell~$2$ at
$d = 2$ (Mukai \texttt{arXiv:math/0408129}~Thm.~0.1); its
$\Gm\times\langle\iota\rangle$-action on the Mukai lattice
$\widetilde H^{\bullet}(K3,\Z)$ twists the chart-wise CoHA by a
$\langle\iota\rangle$-cocycle. Assembling along the quotient
$X/\Aut_{s}(X)$, the cocycle class lives in
$H^{1}(X/\Aut_{s}(X),\underline{\Aut_{s}(X)})$.

\emph{Cell 3 (orbifold alone).} For $X=\C^{3}/G$ with
$G\subset SU(3)$ finite, the inertia stack $I(X/G)$ decomposes by
conjugacy class of $G$. The Bridgeland--King--Reid equivalence
$D^{b}(X)\simeq D^{b}_{G}(\C^{3})$ identifies the global CoHA with a
$G$-equivariantised critical CoHA. The gluing cocycle classifies the
$G$-twists, i.e. $H^{1}(I(X/G),\underline{G})$.

\emph{Cell 4 (automorphic alone).} At a Humbert divisor
$\cH_{D}\subset\cD_{L}$ in the period domain for a lattice $L$ of
signature $(2,n)$, the period map picks up a
$\Gamma_{\mathrm{arith}}\subset O^{+}(L)$-monodromy; the
Borcherds-lifted modular form (a cusp form of weight
$(\mathrm{rk}\,L-2)/2$) assembles local Fourier coefficients into
global CoHA structure constants
(Borcherds~\texttt{arXiv:alg-geom/9609022}~Thm.~10.1;
Gritsenko--Nikulin~\texttt{arXiv:alg-geom/9504006}~Thm.~2.1). The
cocycle is $H^{1}(\cD_{L}/\Gamma_{\mathrm{arith}},
\underline{\Gamma_{\mathrm{arith}}})$.

\paragraph{Pairwise cells.} Write $\cF_{i},\cF_{ii},\cF_{iii},\cF_{iv}$
for the sheaves-of-groups governing each singleton cell:
$\cF_{i}=\underline{\Pic}_{T^{d}}$,
$\cF_{ii}=\underline{\Aut_{s}(X)}$, $\cF_{iii}=\underline{G}$,
$\cF_{iv}=\underline{\Gamma_{\mathrm{arith}}}$. The fibre/semidirect
dichotomy is dictated by the group-theoretic relationship of the two
strata at the overlap: when the two strata share only the diagonal
$\Gm$-scaling of $\Omega_{X}$ and their group actions commute on the
common chart (cells~$5,7,9$), the cocycle sheaf is the fibre product
$\cF_{a}\times_{\underline{\Gm}}\cF_{b}$ over the shared
$\Gm$-diagonal; when one stratum's group acts on the other's
structure by outer automorphism --- the orbifold $G$ permuting toric
chart indices ($\cF_{i}$--$\cF_{iii}$ in cell~$6$), or $\Aut_{s}$
permuting $M_{24}$-conjugacy classes ($\cF_{ii}$--$\cF_{iii}$ in
cell~$8$) --- the cocycle sheaf is the semidirect product
$\cF_{a}\rtimes\cF_{b}$.

\emph{Cell 6 (local $\mathbb{P}^{2}$), derived explicitly.} Local
$\mathbb{P}^{2}=\mathrm{Tot}(K_{\mathbb{P}^{2}})\cong
[\C^{3}/\Z_{3}]^{\mathrm{res}}$: the crepant resolution of
$\C^{3}/\Z_{3}$ with weights $(1,1,1)$. The orbifold stratum
$\Str_{iii}$ sees $G=\Z_{3}$ acting on $\C^{3}$; the toric stratum
$\Str_{i}$ sees the crepant resolution as a fan with three
top-dimensional cones meeting at the exceptional $\mathbb{P}^{2}$.
BKR gives $D^{b}([\C^{3}/\Z_{3}])\simeq D^{b}(\mathrm{local}\,
\mathbb{P}^{2})$; the skew-group algebra
$\C[x,y,z]\rtimes\Z_{3}$ is Morita-equivalent to the Beilinson
endomorphism algebra $\End(\cO\oplus\cO(1)\oplus\cO(2))$ on
$\mathbb{P}^{2}$. The cell-$6$ cocycle $H^{1}(I(X/G),\underline{\Pic}
\rtimes\underline{G})$ concretely classifies the $\Z_{3}$-equivariant
Fourier--Mukai twists: each fan-edge gives a $\Pic$-twist, and the
whole configuration carries a cyclic $\Z_{3}$-action permuting the
three cones. The CoHA transition is the Neguț shuffle-algebra wheel
condition; verifying cell-$6$ against cell-$6$'s alternative chart-wise
presentation (either McKay orbifold or fan-wise toric) gives the
cocycle equivalence.

\emph{Cell 8 (Mathieu K3), derived explicitly.} A Mathieu-K3 is a K3
surface with $G\subset M_{24}$ acting symplectically. The orbifold
stratum gives a $G$-equivariant K3, and the reduced $\Aut_{s}$
stratum gives the $\Aut_{s}(K3)$ action compatible with the Mukai
lattice. The cell-$8$ cocycle
$H^{1}(I(X/G)/\Aut_{s},\underline{\Aut_{s}}\rtimes\underline{G})$
carries the Mathieu moonshine weight-$1/2$ mock modular data: for
each $M_{24}$-conjugacy class $[g]$ there is a twined elliptic genus
$\phi_{g}\in J_{0,1}(\Gamma_{0}(N_{g}))$ (Eguchi--Ooguri--Tachikawa
2011), and the cell-$8$ cocycle assembles the $\phi_{g}$ into a
global CoHA twist. The $\Aut_{s}$-semidirect factor is the action of
$\Aut_{s}(K3)$ on the conjugacy-class indexing set of $M_{24}$
by outer automorphism.

\paragraph{Triple cells.} Cells $11$--$14$ nest three strata; each
admits a concrete witness. The pivotal one for Vol~III is cell $14$:
$K3\times E$ in strata $(ii,iii,iv)$ without a toric input. The
cocycle group is the threefold-nested semidirect product
\[
H^{1}\bigl(\cD_{L}/(\Aut_{s}\times M_{24}),\;
\underline{\Aut_{s}(K3\times E)}\rtimes\underline{M_{24}}\rtimes
\underline{\Gamma_{\mathrm{arith}}}\bigr),
\]
on the lattice-polarised period domain $\cD_{L}$ of $K3\times E$
(with $L$ the transcendental lattice of signature $(2,3)$ giving
rise to the paramodular $\Gamma_{\mathrm{arith}} = K(1) \subset
\Sp_{4}(\Q)$), quotiented by the symplectic automorphism and
Mathieu $M_{24}$-inertia actions.

\paragraph{Cell 15.} For $K3\times E$ with its elliptic fibration
$\pi:K3\to\mathbb{P}^{1}$ providing a partial toric data input (the
fibre at each of the $24$ $I_{1}$-cusps), all four strata meet. The
cocycle group is the fourfold-nested fibre product; no other Vol~III
witness sits at this cell. The $24$-stalk assembly of the $K3\times E$
CoHA is the explicit chain-level manifestation of the cell-15 cocycle.

Elementary derivation: each $I_{1}$-singular fibre of $\pi$ is a
nodal cubic $\{xy=0\}\subset\C^{2}$, and $I_{1}\times E\subset
K3\times E$ is a curve-supported locus. Locally at a node, the CY$_3$
looks like $(\{xy=0\}\times E)^{\mathrm{deformed}}\subset\mathrm{smooth}$,
which is not a $\C^{3}$-chart but a curve-stalk carrying the
Jordan-triple potential $W=\mathrm{tr}(x[y,z])$ on the transverse
directions, tensored with $E$-coefficients. Davison's critical CoHA
on this stalk is computable: it is a curve-stalk version of the
$\C^{3}$-Yangian $Y^{+}(\widehat{\mathfrak{gl}}_{1})$, with the
$E$-direction contributing elliptic parameters to the Mellin-transform
structure functions. Twenty-four such stalks plus the three other
strata ($\Aut_{s}$, Mathieu, Borcherds) assemble via the cell-$15$ cocycle.
The resulting $\kBKM$ is $c_{1}(0)/2=5$ where $c_{1}(0)=10$ is the
zeroth Fourier coefficient of Gritsenko's $\Delta_{5}$; this is the
chain-level source of $\kBKM(\mathfrak{g}_{\Delta_{5}})=5$ in the
K3$\times$E spectrum $\{0,3,5,24\}$.

\paragraph{Cell-by-cell status.} Cells $1$--$4$ (singleton) are
theorem-level per the primary-source citations: cell~$1$ via
Davison--Meinhardt \texttt{arXiv:1601.02479}~Thm.~A for the toric
tilting cocycle; cell~$2$ via Mukai \texttt{arXiv:math/0408129}~
Thm.~0.1 for K3 with Nikulin involution at $d = 2$ and the
translation-equivariant critical CoHA on abelian $A^{3}$ at $d = 3$;
cell~$3$ via BKR \texttt{arXiv:math/9908027}~Thm.~1.1; cell~$4$ via
Borcherds \texttt{arXiv:alg-geom/9609022}~Thm.~10.1 and
Gritsenko--Nikulin \texttt{arXiv:alg-geom/9504006}~Thm.~2.1. Cells
$5$--$11$ and $14$ are instantiated by existing Vol~III
constructions (toric; local $\mathbb{P}^{2}$; Mathieu $K3$;
Borcherds-lifted $K3$; $K3 \times E$ in strata $(ii,iii,iv)$) with
chart-by-chart gluing cocycles verified at the structural level but
not cohomologically computed. Cells $12$ and $13$ lack explicit
Vol~III witnesses: a toric-cum-$\Aut_{s}$-cum-lattice-polarised
CY$_{3}$ (cell $12$) and a toric-cum-orbifold-cum-lattice-polarised
CY$_{3}$ (cell $13$) are each candidate constructions requiring a
simultaneous-three-stratum CY input. Cell~$15$ ($K3 \times E$ with
the elliptic fibration supplying a toric germ at each $I_{1}$-node)
is the deepest instantiated cell.

\paragraph{BCOV anomaly propagation across strata.} Orthogonal to
the cocycle-cohomology classification, the Costello--Li
Atiyah--dilaton cocycle
\[
 \alpha_{\mathrm{BCOV}}(X)
 \;=\; \frac{\chi_{\mathrm{top}}(X)}{24}\,\tr\,\mathrm{At}(T_X)
 \;\in\; H^{1}(X,\cO_X)
\]
(\texttt{arXiv:1606.00365} Prop.~5.2; cross-ref
\S\ref{sec:gluing:obstructions} proof of $(\mathbf{O_2})$) traces a
second axis through the fifteen cells. The target is the
\emph{structure-sheaf} first cohomology, one scalar obstruction per
$X$, not any symmetric power of $T_X^\ast$. The axis factor-splits:
the topological scalar $\chi_{\mathrm{top}}(X)/24$ depends only on
$c_3(T_X)$ and is insensitive to stratum occupancy; the cohomology
factor $\tr\,\mathrm{At}(T_X)$ sees stratum structure through the
Bogomolov--Beauville decomposition of $X$. The BCOV axis propagates
through $\Stk_{\Str}$ by \emph{four} geometric cases, parallel to
$(\beta_1)$--$(\beta_4)$ of \S\ref{sec:gluing:obstructions}:

\emph{$(\beta_1)$ Non-compact retractable, cells containing stratum
$(i)$.} Cells $1, 5, 6, 7, 11, 12, 13, 15$. Bulk
$H^{1}(X_{\mathrm{con}},\cO)=0$ because the affine projection
$\pi:X_{\mathrm{con}}\to\mathrm{affine\ base}$ collapses bulk
Dolbeault cohomology onto the $1$-skeleton of the toric fan. The
bulk Atiyah--dilaton class vanishes; the non-zero topological
scalar $\chi_{\mathrm{top}}(\bY)=2$ (conifold) or
$\chi_{\mathrm{top}}(\mathrm{local}\,\bP^2)=3$ (local $\bP^2$)
survives as a boundary curving, computed via $\zeta$-regularisation
on the Sasaki--Einstein link $L^5$ of $X_{\mathrm{con}}$
(Klebanov--Witten \texttt{arXiv:hep-th/9807080}:
$L^5(\bY)=T^{1,1}\simeq S^2\times S^3$;
$L^5(\mathrm{local}\,\bP^2)=S^5/\Z_3$). The boundary integrand is a
fermionic $\beta\gamma$-ghost system at central charge $c=-2$
(Kausch \texttt{arXiv:hep-th/9510149} Eq.~(4.12); weight-$(1,0)$
$\beta\gamma$), and $\zeta(0)$-regularisation of the link spectrum
fixes the residual curving:
$\kappa_{\mathrm{ch,BV}}(\bY)=-1$,
$\kappa_{\mathrm{ch,BV}}(\mathrm{local}\,\bP^2)=3/2$. The BKR
equivalence (Bridgeland--King--Reid \texttt{math/9908027}) supplies
the skew-group-algebra trivialisation on cell $6$; the local
$\bP^2$ curving is the local instance of the $\zeta$-regularised
ghost-mode residue after McKay resolution.

\emph{$(\beta_2)$ Compact strict CY$_3$ (no cell in $\Stk_{\Str}$,
fall-back).} The quintic $X_5$, the bicubic $X_{3,3}$, the Schoen
threefold. Strict-CY Hodge gives $h^{0,1}(X)=h^{1,0}(X)=0$, so the
cohomology target $H^1(X,\cO_X)$ is itself zero. The
Atiyah--dilaton class is zero trivially; genuine obstructions lie
in higher BCOV weights via iterated Atiyah-class towers outside the
leading cocycle, and outside the $\Stk_{\Str}$-cell classification
(these CY$_3$ admit no stratum assignment). Cross-ref
Lemma~\ref{lem:chi-top-bcov-split} of
\S\ref{sec:gluing:obstructions}.

\emph{$(\beta_3)$ $K3\times E$ at cell $15$.} Künneth collapses the
topological scalar: $\chi_{\mathrm{top}}(K3\times E)=\chi(K3)\cdot
\chi(E)=24\cdot 0=0$. Independently,
$h^{0,1}(K3\times E)=h^{0,1}(E)=1$, so $\tr\,\mathrm{At}(T_X)\neq
0$ in $H^1(K3\times E,\cO)$. The product
$\alpha_{\mathrm{BCOV}}(K3\times E)=0$ vanishes for
\emph{topological} reasons. The $\kappa_{\mathrm{ch,BV}}$ curving
on the cell-$15$ stalk is supplied by the stratum-$(iv)$ Gritsenko
$\Delta_5$-Borcherds lift, not by a boundary-link $\zeta$-residue;
it aligns with the Heisenberg--Mukai invariant
$\kappa_{\mathrm{ch}}^{\mathrm{Heis}}(K3\times E)=3$ on the
stratum-$(iv)$ side.

\emph{$(\beta_4)$ Compact CY$_3$ with abelian factor, cells
containing $\Str_{ii}$ or $\Str_{iv}$.} Cells $2, 4, 9, 10, 14$.
The elliptic-curve cell $2$, the abelian-factor lattice-polarised
cell $4$, the pairwise and triple cells $9, 10, 14$ containing an
abelian summand in Bogomolov--Beauville.
$\chi_{\mathrm{top}}=0$ and $h^{0,1}\neq 0$; the cocycle vanishes
topologically while the cohomology target is non-trivial, so the
deformation \emph{gradient} of $\alpha_{\mathrm{BCOV}}$ in the cell
detects the abelian direction rather than the cocycle class
itself. The purely-orbifold cell $3$ ($[\C^3/G]$ with $G\subset
\mathrm{SU}(3)$ finite, no abelian factor) sits in $(\beta_1)$:
retraction onto the exceptional divisor collapses bulk
$H^1(\widetilde{X},\cO)=0$, and the boundary $\zeta$-residue
supplies $\kappa_{\mathrm{ch,BV}}$.

\emph{Ghost-mode balance on class-G cells.} The class-G free-field
subfamily of cell $1$ (resolved conifold, SPP, banana --- Polyakov
cone geometries with free-field boundary-link treatment) satisfies
\[
 \kappa_{\mathrm{ch}}(X) + \kappa_{\mathrm{ch,BV}}(X)
 \;=\; \kappa_{\mathrm{cat}}(X)
\]
at $\chi_{\mathrm{top}}=2$, realised on the conifold as the
canonical four-invariant row
$(\kappa_{\mathrm{ch}},\kappa_{\mathrm{cat}},
\kappa_{\mathrm{BKM}},\kappa_{\mathrm{ch,BV}})(\bY)=
(+1,0,+1,-1)$ of
Theorem~\ref{thm:conifold-polyakov-ghost-balance} in
\texttt{cy\_d\_kappa\_stratification.tex}. The balance is
\emph{not} the Theorem~C derived-centre ceiling
$K^\kappa\in\{0,8,13,250/3,98/3\}$ and is not universal across
$\Stk_{\Str}$; it fails on cells outside the class-G free-field
cone family.

\emph{The (O2) contribution to $\Stk_{\Str}$.} The BCOV axis is
independent of the cocycle classification of \S 10.1: the cocycle
group $H^{1}(\cdot,-)$ at each cell addresses assembly of critical
CoHAs from local data, whereas $\alpha_{\mathrm{BCOV}}\in
H^{1}(X,\cO_X)$ is the Hodge-theoretic Atiyah--dilaton class
governing quantisation. The two coexist at every cell: a cell with
vanishing $\alpha_{\mathrm{BCOV}}$ but non-trivial cocycle (toric
CY without compact $4$-cycles at cell $1$) has an assemblable CoHA
with $\zeta$-regularised boundary curving; a cell with non-zero
$\alpha_{\mathrm{BCOV}}$ but trivial cocycle (hypothetical rigid
strict CY$_3$ in no stratum of $\Stk_{\Str}$) has neither assembly
nor bulk quantisation. The product cell $K3\times E$ at cell $15$
is the case where both meet: the cocycle is the fourfold nested
fibre product, $\alpha_{\mathrm{BCOV}}(K3\times E)=0$ by Künneth
despite $\tr\,\mathrm{At}(T_X)\neq 0$, and the
$\kappa_{\mathrm{ch,BV}}$ curving coincides with the
Heisenberg--Mukai invariants on the stratum-$(iv)$ side.

\emph{Cross-volume (Vol II 3D HT QFT).} The stratum-selective BCOV
behaviour on $\Stk_{\Str}$ is the Vol~III shadow of Vol~II's
BCOV-anomaly obstruction to $E_{3}$-chiral (Vol~II main.tex~line
$1542$: $\kappa_{\mathrm{BCOV}} = \chi(X)/24$ for Class B CY outside
the K3-fibered locus). Vol~II records that the BCOV invariant is the
obstruction to the chiral-side lift being $E_{3}$-chiral rather than
merely $E_{1}$-chiral; Vol~III~\S 10 refines this by stratum: the
cells $1$--$15$ carrying $\alpha_{\mathrm{BCOV}} = 0$ admit the
$E_{3}$-chiral lift through the Stage-1 $\Phi^{FA}_{3}$ (the
factorisation algebra is quantisable), and the cells with
$\alpha_{\mathrm{BCOV}} \neq 0$ drop to $E_{1}$-chiral after
Stage-2 specialisation. The three-volume agreement: Vol~II states
the obstruction abstractly, Vol~III stratifies the vanishing locus
across cells $1$--$15$, Vol~I tracks the bar-cobar shadow on the
stratum-resolved chiral lift.

\subsubsection{\S 10.2. The sheaf of CoHAs on the $(2,1)$-stack of
equivariance strata}
\label{gluing:sec:10-2-sheaf-on-stack}

The classification of \S 10.1 organises CoHAs by the cell of a
combinatorial structure (the poset of fifteen cells); we upgrade this
to a derived-geometric statement by exhibiting the cells as a
$(2,1)$-stack. The terminology follows Toen--Vezzosi and Lurie.

\begin{definition}[\ClaimStatusDefinition\ The stratum stack]
\label{gluing:def:stratum-stack}
$\Stk_{\Str}$ is the $(2,1)$-stack over the \'etale site of derived
Artin stacks locally of finite presentation (Toen--Vezzosi
\texttt{arXiv:math/0404373}~\S 1.3; Lurie HTT~\S 5.1) defined by:
\begin{itemize}[leftmargin=1.5em]
 \item \textbf{objects}: pairs $(X,S)$ with $X$ a CY$_{d}$ derived
 Artin stack and $S\subseteq\{i,ii,iii,iv\}$ a choice of strata with
 $X$ carrying each stratum datum indexed by $S$ (a toric fan when
 $i\in S$; a reduced $\Gm\times\Aut_{s}(X)$-action when $ii\in S$; an
 orbifold $[Y/G]$-presentation with $G\subset SU(d)$ finite when
 $iii\in S$; a period map $X\to\cD_{L}/\Gamma_{\mathrm{arith}}$ to
 a lattice-polarised domain intersecting a Humbert-type special
 divisor $\cH_{D}$ when $iv\in S$);
 \item \textbf{$1$-morphisms} $(X,S)\to(X',S')$: derived-stack maps
 $f:X\to X'$ with $S\supseteq S'$ and each stratum datum of $S'$ the
 $f$-pullback of the corresponding datum in $S$;
 \item \textbf{$2$-morphisms}: natural isomorphisms of
 stratum-datum-preserving maps, i.e.\ coherent homotopies between
 $1$-morphisms through the stratum-datum structure.
\end{itemize}
The topology is fppf. The sheaf condition is the homotopy-coherent
version of fppf-descent for stratum-datum-equipped CY data
(Toen--Vezzosi \emph{loc.~cit.}~Thm.~1.3.7.2 for the derived-stack
$\Cech$ descent).
\end{definition}

\begin{remark}[First-principles reading]
A point of $\Stk_{\Str}$ is a CY $X$ with a declared stratum occupancy
$\Str(X)\subseteq\{i,ii,iii,iv\}$. A $1$-morphism
$(X,S)\to(X',S')$ restricts to a sub-stratum occupancy $S\supseteq S'$
and carries the extra datum of the stratum-compatibility of
$f:X\to X'$. A $2$-morphism is a homotopy between two such
$1$-morphisms, coherent with the stratum-compatibility data.
\end{remark}

\begin{theorem}[\ClaimStatusConjectured\ The sheaf-of-CoHAs]
\label{gluing:thm:sheaf-of-cohas}
There exists a sheaf $\underline{\cH}$ on $\Stk_{\Str}$ whose stalk at
a CY $X\in\Stk_{\Str}$ is the global CoHA $\cH(X)$, and whose
transition cocycles on overlaps between stratum cells are precisely the
gluing cocycles of \S 10.1.
\end{theorem}

\begin{proof}[First-principles construction, four steps]
\emph{Step~1 (local model).} For $X$ in a singleton cell of
$\Stk_{\Str}$, the local CoHA is the Davison--Meinhardt critical CoHA
(\texttt{arXiv:1601.02479}~Theorem~5.11: the vanishing-cycle sheaf
$\phi_{W}\cO_{\cM^{ss}_{X}}$ on the semistable moduli stack
$\cM^{ss}_{X}$ carries a Hall product from the extension-closed
substack correspondence): the Jacobi algebra of a $\C^{3}$-chart
(cell~$1$, \texttt{arXiv:1202.2756}~Theorem~1.1); the
$\Aut_{s}(X)$-equivariant critical CoHA on $X/\Aut_{s}$ (cell~$2$);
the $G$-equivariantised critical CoHA on $[\C^{3}/G]$ via
Bridgeland--King--Reid (cell~$3$,
\texttt{arXiv:math/9908027}~Theorem~1.1); the Borcherds-lifted
Fourier-coefficient assembly on the period domain (cell~$4$,
\texttt{arXiv:alg-geom/9609022}~Theorem~10.1). Each is a well-defined
stable $\infty$-categorical algebra object (Lurie HA~\S 7.3).

\emph{Step~2 (descent along $1$-morphisms).} A $1$-morphism
$f:X\to X'$ in $\Stk_{\Str}$ is a stratum-preserving CY map. The
pullback $f^{*}\phi_{W'}\cO_{\cM^{ss}_{X'}}$ of the vanishing-cycle
sheaf agrees with $\phi_{W}\cO_{\cM^{ss}_{X}}$ on strata-preserving
inclusions by Davison \texttt{arXiv:1311.6989}~Proposition~3.13
(critical CoHA pullback along flat maps commutes with the Hall
product). Hence the local CoHAs of Step~$1$ assemble into a presheaf
of stable $\infty$-categorical algebras on the $1$-category underlying
$\Stk_{\Str}$.

\emph{Step~3 (coherence via $2$-morphisms).} A $2$-morphism in
$\Stk_{\Str}$ is a homotopy between stratum-shift choices (the
$(2,1)$-truncation records homotopies; see
Remark~\ref{gluing:rem:2-1-truncation-justification} below). The
associated gluing of local CoHAs lands in the cocycle group named in
the table of \S 10.1: for pairwise cell $(ii,iii)$ one gets
$H^{1}(I(X/G)/\Aut_{s},\underline{\Aut_{s}}\rtimes\underline{G})$
(Mathieu twist combined with reduced $\Aut_{s}$-cocycle); for
triple/quadruple cells the cocycle is a nested fibre or semidirect
product of the four basic sheaves of groups.

\emph{Step~4 (\v{C}ech-sheaf assembly).} Given a cover of $\Stk_{\Str}$
by stratum-cell opens, the cell-local CoHAs glue along the Step~$3$
cocycle to produce a sheaf on $\Stk_{\Str}$. The sheaf axioms are the
homotopy-coherent analogue of \v{C}ech descent for sheaves of stable
$\infty$-categories (Toen--Vezzosi \texttt{arXiv:math/0404373}~\S 1.4;
Lurie HTT~\S 6.2.3). Sections over a CY $X\in\Stk_{\Str}$ recover
$\cH(X)$.
\end{proof}

\begin{remark}[Why $(2,1)$ and not $(1,1)$, $(2,2)$, or $(\infty,1)$]
\label{gluing:rem:2-1-truncation-justification}
The stratum-inclusion poset on a fixed CY $X$ is the power set
$\cP(\{i,ii,iii,iv\})\setminus\{\varnothing\}$, a finite incidence
poset of length $4$. This constrains the truncation as follows.

\emph{Not $(1,1)$.} Two choices of stratum-filtration on the same $X$
can differ by an element of $\Aut_{s}(X)\times G\times
\Gamma_{\mathrm{arith}}$ acting as an isomorphism within the
stratum-datum category. A $(1,1)$-stack identifies any two such
choices, erasing the $\Aut_{s}$-inner-automorphism twists that
distinguish, e.g., cells~$8$ and~$10$ on the same Mathieu~K3. The
semidirect structure $\underline{\Aut_{s}}\rtimes\underline{G}$ is
lost without a $2$-morphism layer recording the Aut-conjugation
data.

\emph{Not $(2,2)$.} All $2$-morphisms in the stratum-datum category
are natural isomorphisms: a homotopy between two stratum-preserving
maps $f_{0},f_{1}\colon (X,S)\to(X',S')$ is a coherent invertible
transformation through the stratum-compatibility, there being no
non-invertible $2$-morphism structure on the datum itself. The
$(2,1)$-truncation retains only invertible $2$-morphisms, which is
exactly what the stratum-compatibility furnishes.

\emph{Not $(\infty,1)$.} A full $(\infty,1)$-stack would record
$n$-morphisms for $n\geq 3$. These are not load-bearing: the
stratum-inclusion poset has length~$4$, so its nerve has dimension
$\leq 3$ as a simplicial complex, and the space of stratum-filtrations
on a fixed $X$ is discrete up to a finite-dimensional $\Aut_{s}\times
G\times\Gamma_{\mathrm{arith}}$-action whose orbit spaces are
classifying stacks $B\Aut_{s}$, $BG$, $B\Gamma_{\mathrm{arith}}$
--- already $(2,1)$-truncated. Francis--Gaitsgory
\texttt{arXiv:1103.5803} Thm.~1.2.4 + 5.2.1 on factorisation-homology
descent for finite-incidence diagrams gives the explicit bound: for
a finite poset $\cP$ of length $\ell$, factorisation descent on the
associated nerve lands in the $(\ell-1,1)$-truncation; at $\ell=4$
this is $(3,1)$, and the further restriction to stratum-compatibility
(which is $2$-categorical by the Aut-semidirect argument above) cuts
this to $(2,1)$.

Thus level $0$~$=$~CY data, level $1$~$=$~stratum-preserving
morphisms, level $2$~$=$~coherence between stratum-compatibility
witnesses, and levels $\geq 3$ contract. The $(2,1)$-truncation is
sharp: coarser truncations lose load-bearing semidirect data, finer
truncations record data that does not exist.
\end{remark}

\paragraph{The $K3\times E$ stalk at cell-$15$ (chain-level).} The
cell-$15$ point of $\Stk_{\Str}$ hosting $K3\times E$ with its
elliptic fibration $\pi:K3\to\mathbb{P}^{1}$ has stalk
\[
\underline{\cH}_{K3\times E}
\;=\;
\mathrm{hocolim}_{\Delta^{3}_{\Str}}
\bigl(\cH^{(i)}_{\mathrm{curve}},\;\cH^{(ii)}_{\Aut_{s}},\;
\cH^{(iii)}_{M_{24}},\;\cH^{(iv)}_{\Delta_{5}}\bigr),
\]
where $\Delta^{3}_{\Str}$ is the simplicial nerve of the poset
$\cP(\{i,ii,iii,iv\})\setminus\{\varnothing\}$ restricted to chains
terminating at $\{i,ii,iii,iv\}$, equivalently the $3$-simplex whose
vertices are the four singleton cells and whose higher faces are the
pairwise, triple, and quadruple overlaps (Lurie HTT~\S 2.3.4 on
simplicial nerves of finite posets),
where each of the four inputs is chain-level explicit by
Theorem~\ref{gluing:thm:k3xe-four-stratum} of \S\ref{sec:gluing:k3e}:
\begin{enumerate}[label=\textup{(\alph*)}]
 \item $\cH^{(i)}_{\mathrm{curve}}$: twenty-four $I_{1}\times E$
 curve-stalks indexed by the $24$ singular fibres of $\pi$;
 Picard--Lefschetz monodromy $T_{b}\in\SL_{2}(\Z)$ twists the Hall
 product on each (Proposition~\ref{gluing:prop:local-model-curve-stalk}).
 At the generic K3-point each curve-stalk computes to the
 Schiffmann--Vasserot algebra $\Yplus(\widehat{\fgl}_{1})\otimes H^{*}(E)$
 (\texttt{arXiv:1202.2756} Theorem~1.1, base-changed to $E$).
 \item $\cH^{(ii)}_{\Aut_{s}}$: the reduced
 $\Aut_{s}(K3)\times\Aut(E)$-cocycle on
 $X/(\Aut_{s}(K3)\times\Aut(E))$; $\Aut(E)=E\rtimes\Z/2$ with $E$
 acting by translation and $\Z/2$ by inversion.
 \item $\cH^{(iii)}_{M_{24}}$: the Mathieu moonshine inertia
 $I(K3/M_{24})\to M_{24}\mathrm{\text{-}Conj}$; each conjugacy class
 $[g]$ contributes the twined elliptic genus
 $\phi_{[g]}\in J_{0,1}(\Gamma_{0}(N_{[g]}))$
 (\texttt{arXiv:1004.0956}~\S 3).
 \item $\cH^{(iv)}_{\Delta_{5}}$: the Borcherds form $\Delta_{5}$ on
 the Kuga--Satake period domain, weight $5$, Fourier coefficient
 $c_{1}(0)=10$; the gluing cocycle is the automorphic-form cocycle
 on $\cA_{2}$ restricted to the period image of $K3\times E$
 (Gritsenko--Nikulin \texttt{arXiv:alg-geom/9504006} Pt.~II Thm.~2.1;
 Borcherds \texttt{arXiv:alg-geom/9609022}).
\end{enumerate}
The cell-$15$ cocycle is the fourfold nested fibre product of (a)--(d)
over the appropriate intersections of stratum sheaves; its
$\kBKM$-invariant is $\kBKM(\mathfrak{g}_{\Delta_{5}})
=c_{1}(0)/2=10/2=5$, reproducing the cell-$(iv)$ entry
of the $K3\times E$ spectrum $\{0,3,5,24\}$.

\paragraph{Compact CY$_{3}$ caveat: the stratum-free fall-back.} For
compact CY$_{3}$ outside the $K3\times E$ family (quintic, bicubic in
$\mathbb{P}^{2}\times\mathbb{P}^{2}$, Schoen threefold) the point
$X\in\Stk_{\Str}$ occupies no singleton stratum cell: there is no
global toric structure, no residual $\Aut_{s}^{0}$, no finite
orbifold cover, no Humbert-lifted period structure. The sheaf
$\underline{\cH}$ of Theorem~\ref{gluing:thm:sheaf-of-cohas} does not
assemble at such $X$ via the cell-cocycle presentation, but
$\cH(X)$ itself is non-zero and computable directly: it is the
Davison--Meinhardt critical CoHA
\[
\cH(X) \;=\; \CoHA_{\mathrm{crit}}(X) \;=\;
H^{\bullet}_{\mathrm{crit}}(\cM^{ss}_{X},\phi_{W})
\]
(Davison--Meinhardt \texttt{arXiv:1601.02479}~Theorem~5.11) together
with its MNOP categorification at the chain level
(Maulik--Nekrasov--Okounkov--Pandharipande 2006
\texttt{arXiv:math/0312059}~Conjecture~3) and Davison integrality
(\texttt{arXiv:1311.6989}~Theorem~A) providing the PBW generators
of the BPS Lie algebra. Stratum-free content: the sheaf
$\underline{\cH}$ extends to such $X$ as the \emph{constant sheaf}
with stalk $\CoHA_{\mathrm{crit}}(X)$ --- no transition cocycles
between charts of $X$ are required, because the relevant moduli-stack
cohomology $H^{\bullet}_{\mathrm{crit}}$ is already assembled
globally by the critical-cohomology formalism. The stratum-free
fall-back is the codomain of \S\ref{sec:gluing:obstructions}'s
obstruction analysis, providing the Vol~III companion to the
cell-structured cases.

\subsubsection{\S 10.3. Cross-volume relevance}
\label{gluing:sec:10-3-cross-volume}

The sheaf-of-CoHAs-on-$\Stk_{\Str}$ perspective interlocks with the
framework of the prior two volumes of this programme.

The binding datum is the two-stage factorisation of the
CY-to-chiral functor (Vol~III, \texttt{cy\_to\_chiral.tex} Theorem
\texttt{phi-two-stage-factorisation}):
\[
\Phi_{3}
\;=\;
\SpCh_{\Sigma_{2},C}\circ\PhiFA_{3}
\;:\;
\mathrm{CY}_{3}\text{-Cat}\;\longrightarrow\;\EnHolFA(X)\;\xrightarrow{\;\SpCh\;}\;\Eone\text{-ChirAlg}(C).
\]
Stage~$1$ $\PhiFA_{3}$ is canonical up to contractible choice
(Kontsevich--Tamarkin formality + Costello--Gwilliam--Li locality);
Stage~$2$ $\SpCh_{\Sigma_{2},C}$ is specialisation along a $2$-cycle
$\Sigma_{2}\subset X$ restricted to the reference curve $C\subset X$.

\paragraph{Vol~I (chiral bar--cobar, operadic engine).} The
bar--cobar adjunction $(B,\Omega)$ of Vol~I Theorem~A on
$\Eone$-chiral algebras operates \emph{after} Stage~$2$: given a CY
datum $X\in\Stk_{\Str}$ with specialisation input $(\Sigma_{2},C)$,
the $\Eone$-chiral algebra $\Phi_{3}(X)$ admits a bar complex
$B(\Phi_{3}(X))$ whose chiral cobar $\Omega B(\Phi_{3}(X))\simeq
\Phi_{3}(X)$ by Vol~I Theorem~A. The $\Stk_{\Str}$-cocycle class on
$X$ pulls back through $\Phi_{3}$ to a cocycle on
$B(\Phi_{3}(X))$: the Vol~I bar--cobar is the shadow, on the
reference curve $C$, of the $\Stk_{\Str}$-gluing datum on $X$. Bar--
cobar is the operadic engine; the $\Stk_{\Str}$-cocycle is the
geometric one; they agree after $(\SpCh,\Phi)$-transport.

\paragraph{Vol~II (3D HT QFT, geometric engine).} The 3D
holomorphic-topological twisted partition function of Vol~II decomposes
the 3-manifold $M^{3}=\Sigma_{2}\times C$ via factorisation-homology
gluing: the $E_{3}$-factorisation algebra of Stage~$1$, $\PhiFA_{3}(X)
\in\EnHolFA(X)$, integrates over $\Sigma_{2}\subset X$ via
$\SpCh_{\Sigma_{2},C}=\int_{\Sigma_{2}}(-)|_{C}$ (Dunn--Lurie
decomposition $E_{3}\simeq E_{2}\otimes E_{1}$, Lurie HA~\S 5.1.2;
Costello--Gwilliam \texttt{arXiv:2210.13036}~\S 4 for stratified
factorisation homology). Vol~II thus furnishes the geometric engine
on $M^{3}=\Sigma_{2}\times C$: Stage~$1$ lives on $X$, the transverse
$\Sigma_{2}$-integration lands on $C$, and the resulting
$\Eone$-chiral algebra on $C$ is the $\Phi_{3}$-image ready for
bar--cobar.

\paragraph{Vol~III (CY quantum groups, combinatorial--cellular engine).}
Vol~III contributes \S 10.1--10.2 above: the classification of CoHAs
by stratum cell and the sheaf-of-CoHAs presentation. The $K3\times E$
CoHA is the cell-$15$ master witness; the Mathieu $K3$ CoHA is
cell-$8$ or cell-$10$ (depending on Humbert specialisation); the
local $\mathbb{P}^{2}$ CoHA is cell-$6$; the $\C^{3}$ CoHA is
cell-$1$.

\paragraph{Three-engine coherence.} The three engines descend from a
single cocycle in $\Stk_{\Str}$ via three distinct pushforwards
through $\PhiFA_{3}$:
\[
[c_{X}]\in H^{1}(X,\underline{F}_{\Str(X)})
\quad\xrightarrow{\;\PhiFA_{3}\;}\quad
\begin{cases}
\text{bar}_{C}\text{: Vol~I cocycle on }B(\Phi_{3}(X)|_{C}),\\[0.5ex]
\int_{\Sigma_{2}}\text{: Vol~II cocycle on }M^{3}=\Sigma_{2}\times C,\\[0.5ex]
\mathrm{id}_{X}\text{: Vol~III cocycle on }\cH(X).
\end{cases}
\]
The coherence of the three images is the Dunn--Lurie decomposition
$E_{3}\simeq E_{2}\otimes E_{1}$ (Lurie HA \S 5.1.2.4): the Vol~II
$\Sigma_{2}$-integration extracts the $E_{2}$-factor, bar--cobar on
$C$ extracts the $E_{1}$-factor, and Vol~III's CoHA is the total
$E_{3}$-algebra on $X$. On singleton cells the agreement is direct
(single stratum sheaf throughout). On pairwise and higher cells the
agreement passes through the nested fibre/semidirect structure of
the \S 10.1 cocycle, which Dunn--Lurie factorises stratum-by-stratum
along $E_{3}\simeq E_{2}\otimes E_{1}$.

\subsubsection{\S 10.4. Open frontiers}
\label{gluing:sec:10-4-open-frontiers}

Four concrete frontiers remain beyond the present framework.

\paragraph{(F1) Verification of $\Stk_{\Str}$'s descent axiom.}
Definition~\ref{gluing:def:stratum-stack} specifies
$\Stk_{\Str}$ as a $(2,1)$-stack over the fppf site of derived Artin
stacks locally of finite presentation; the descent axiom is the
fppf-$\Cech$-descent for stratum-datum-equipped CY data, asserted on
the authority of Toen--Vezzosi \texttt{arXiv:math/0404373}~Thm.~1.3.7.2.
The verification that fppf-descent lifts from the underlying derived
stack to the stratum-datum level requires a separate calculation for
each of the four stratum data: toric-fan descent (straightforward;
fans are \'etale-local), $\Aut_{s}(X)$-equivariance descent
(non-trivial under central extensions), $G$-orbifold descent
(BKR-controlled), Humbert-polarisation descent (controlled by
Heegner-divisor geometry, non-trivial). Sheaves of stable $\infty$-categories
on a $(2,1)$-stack are Lurie HTT~\S 6.2.3; the $(2,1)$-truncation
suffices by Remark~\ref{gluing:rem:2-1-truncation-justification}.

\paragraph{(F2) Cocycle cohomology in each of the 15 cells.} For each
of the fifteen cells, \S 10.1 named a cohomology group in which the
gluing cocycle lives, but in no cell beyond~$1$ is that cohomology
explicitly computed. For cell $14$ ($K3 \times E$ in strata
$(ii,iii,iv)$) the cocycle is
\[
H^{1}\bigl(\cD_{L}/(\Aut_{s}\times M_{24}),\;
\underline{\Aut_{s}(K3\times E)}\rtimes\underline{M_{24}}\rtimes
\underline{K(1)}\bigr)
\]
with $K(1) \subset \Sp_{4}(\Q)$ the paramodular group stabilising
$\Delta_{5}$; making this computation concrete requires identifying
the three-factor semidirect structure explicitly (the $M_{24}$-action
on $\Aut_{s}$ by outer automorphism plus the paramodular-action on
the $M_{24}$-conjugacy indexing), parametrising the site, and
computing the first cohomology. For cell $15$ the computation stacks
once more with the $\underline{\Pic}_{T^{1}}$-factor from the
elliptic-fibration toric germ. Each cell is a self-contained
cohomology computation.

\paragraph{(F3) 24-stalk assembly of $K3\times E$ at explicit chain
level.} The $K3\times E$ CoHA at cell $15$ assembles from $24$
curve-stalks on $I_{1}\times E$ (one per singular fibre of the
elliptic fibration $\pi:K3\to\mathbb{P}^{1}$), plus the three
non-toric stratum inputs: $\cH^{(ii)}_{\Aut_{s}}$ from
$\Aut_{s}(K3\times E)$, $\cH^{(iii)}_{M_{24}}$ from the Mathieu
conjugacy data, $\cH^{(iv)}_{\Delta_{5}}$ from the Borcherds lift.
At the level of the critical CoHA on moduli stacks, this is a
definite prescription; at the level of explicit chain-level
Hall-product formulae, however, it is a multi-step computation:
each $I_{1}\times E$-node is a curve-stalk Jordan-triple potential
tensored with $E$-coefficients, and the $24$ contributions sum
through a Mukai-lattice-valued gluing cocycle. No chain-level
verification of the full $24$-stalk sum yet exists in Vol~III; it
is the outstanding concrete frontier.

\paragraph{(F4) Extension to CY$_{5}$ and Fake Monster / rank-25.}
The classification of \S 10.1 was formulated for CY$_{3}$; the stack
$\Stk_{\Str}$ of \S 10.2 likewise. At $d=5$ on $K3\times K3\times E$,
the Fake Monster Lie algebra of Borcherds 1990 lives, with rank-$25$
Cartan on $\mathrm{II}_{25,1}=\Lambda_{24}\oplus U$. The analogous
stratum stack at $d=5$ requires an expanded family of strata
accommodating the Dunn--Lurie $E_{5}=E_{2}\otimes E_{2}\otimes E_{1}$
decomposition. The question of whether the four-stratum taxonomy of
CY$_{3}$ lifts to a four-plus-new-strata taxonomy at CY$_{5}$ is
open; the Fake Monster cell presumably sits at an analogue of
cell-$15$ in an expanded stack.

\paragraph{Status summary.} Theorems~
\ref{gluing:thm:fifteen-cell-classification} and
\ref{gluing:thm:sheaf-of-cohas} place every Vol~III CoHA construction
at a stratum cell with cocycle valued in a named sheaf of groups.
Cells $1$--$4$: theorem-level at the singleton-cocycle level per the
primary-source citations above. Cells $5$--$11$ and $14$: Vol~III
witnesses with structurally verified gluing data. Cells $12$, $13$:
no explicit witness. Cell $15$: the deepest instantiated cell
($K3 \times E$ with its elliptic fibration supplying the stratum-$(i)$
toric germs at the $24$ $I_{1}$-cusps). Cocycle cohomology is
computed only at cell~$1$ (toric-fan $\underline{\Pic}_{T^{d}}$);
(F1)--(F4) name the outstanding tasks.

\paragraph{$\Phi$ as a sheaf map on $\Stk_{\Str}$.} The CY-to-chiral
functor $\Phi:\mathrm{CY\text{-}cat}\to\ChirAlg$ acts on
$\Stk_{\Str}$ as a sheaf map: for $X$ in cell $c$, the chiral
algebra $\Phi(X)$ carries gluing cocycle
$\Phi_{*}([c_{X}])\in\Phi_{*}H^{1}(X,\underline{F}_{\Str(X)})$,
computed through the two-stage factorisation
$\Phi_{3}=\SpCh_{\Sigma_{2},C}\circ\PhiFA_{3}$ with
$(\Sigma_{2},C)$ the cell-compatible specialisation datum. The seven
faces of $r_{\mathrm{CY}}$ of \S\ref{sec:gluing:wallcrossing} are
seven stratum-cell trajectories through $\Stk_{\Str}$; the full
$r_{\mathrm{CY}}$ theory is the monodromy of the $\Phi$-sheaf along
closed loops in $\Stk_{\Str}$.